\section{大疆的经验：从飞控到产品级可靠性}

\subsection{大疆飞控系统的演进}

杨硕从2013年研究生期间开始参与大疆飞控系统开发。他的导师李泽强教授的许多学生都与大疆有深度合作，这为杨硕提供了从学术到工业的自然过渡。

在大疆的五年间（2013--2018），杨硕经历了飞控系统从基础到成熟的完整演进：
\begin{itemize}
    \item 参与并主导了大疆最早期飞控系统的算法重构
    \item 涉及底层控制算法与软件系统架构的重新设计
    \item 参与了多个消费级产品的开发周期
    \item 后期负责 RoboMaster 机器人比赛项目
\end{itemize}

\subsection{对标波音的可靠性追求}

杨硕分享了一个令人印象深刻的细节：2016年，大疆内部有一次会议，白板上列出了竞争对手，其中最大的名字是\textbf{波音}。当时 Phantom 4 的故障率约为两百万分之一，而波音的故障率可以做到五百万分之一。大疆的目标就是向航空航天级别的系统可靠性看齐。

这种追求的背后逻辑是：当你卖出100万甚至200万台会飞的设备，即使两百万分之一的故障率，在绝对数量上仍然会产生不少事故。因此必须持续提升可靠性。

\begin{importantbox}{从无人机到机器人的可靠性启示}
杨硕指出，大疆之所以提出对标波音的可靠性目标，是因为 Phantom 1 在2013年上市后形成了大规模销售，大量用户使用后开始出现事故，倒逼公司将可靠性提升到航天航空级别。这一经验对人形机器人行业有直接借鉴意义——当未来消费级机器人实现大规模量产时，同样会面临类似的可靠性挑战。当前的机器人公司，连两百分之一甚至两千分之一的故障率都很难做到。
\end{importantbox}

\subsection{大疆的护城河：细节的力量}

杨硕认为大疆在精灵（Phantom）发布12年后，其技术护城河的核心已经不是某项单一技术突破，而是\textbf{各种细节的极致追求}。无人机飞控怎么做、相机怎么调、图传怎么设计，这些在行业内早已有成熟方案，但各家"拼到一块就是不如大疆的产品好"。

这为机器人行业提供了一个重要启示：技术趋同是必然的——成熟技术一定能被大规模复制，但细节堆积起来会产生巨大差异。机器人比无人机的复杂度高出不止一个数量级（四个电机 vs. 数十个关节），因此在细节上的差距只会更加放大。

\subsection{产品生命周期的完整链条}

杨硕强调，一个硬件产品的成功远不止于技术本身，而是需要覆盖完整的产品生命周期：

\begin{enumerate}
    \item \textbf{技术实现}：把基础技术做扎实（大多数创业公司目前集中在这一步）
    \item \textbf{量产}：从原型到规模化生产
    \item \textbf{售后}：用户拿到产品后的故障处理
    \item \textbf{返修}：产品使用一段时间后的性能下降与维修体系
\end{enumerate}

杨硕指出，Tesla 之所以在机器人可靠性方面最接近波音水平，正是因为其机器人业务继承了汽车业务线已经建立起来的完整基础设施——从部署、硬件、数据收集到测试，都有非常完善的体系和流程。

\subsection{本章小结}

大疆的经历为杨硕提供了从学术到工业的完整视角。对标波音的可靠性追求、对细节的极致关注、以及对产品全生命周期的深刻理解，构成了他后来创业的核心方法论。一个核心洞见是：技术护城河不在于某项单一突破，而在于无数细节的持续优化与系统化积累。


